\section{The Coding Scheme Problem}
\label{sec:coding}

\subsection{MDL Requires a Coding Scheme}

The MDL objective is:
\[
  \text{minimize} \quad L(H) + L(D \mid H)
\]
This formula is only well-defined once you specify how hypotheses and
data are encoded---what ``description length'' means. Different coding
schemes give different numerical values of $L(H)$ and $L(D \mid H)$,
and in principle different rankings of theories.

The Kolmogorov framework provides the theoretically correct answer:
the description length of a hypothesis is $\K(H)$, the length of the
shortest program that generates it. The Solomonoff prior
$P(H) \propto 2^{-\K(H)}$ is the resulting prior. This coding scheme
is universal: it dominates every computable distribution up to a
multiplicative constant (Section~\ref{sec:kolmogorov}). For large
enough data, any two universal coding schemes give the same ranking
of hypotheses---the constant washes out.

The problem is that $\K(H)$ is uncomputable. You cannot implement the
Kolmogorov coding scheme directly. Every practical MDL application
uses a restricted coding scheme: a parametric family, a grammar, a
neural architecture. These are computable but not universal---they
cover a restricted hypothesis class and give no principled way to
compare hypotheses from different families.

This is the coding scheme problem: to have a complete formal theory of
science, we need a single coding scheme that is simultaneously
\emph{universal} (covering all scientific hypotheses),
\emph{generative} (assigning likelihoods to data), and
\emph{tractable} (computable in practice).

\subsection{What a Universal Practical Coding Scheme Requires}

A coding scheme for scientific hypotheses must satisfy four
properties to serve as a universal practical approximation to the
Kolmogorov prior:

\begin{enumerate}

\item \textbf{Universality.} The hypothesis space must cover all
scientifically expressible claims. Since all of mathematics and
science can be expressed in first-order logic
(FOL)~\cite{coppola2024qbbn}, a coding scheme is universal in the
relevant sense if it can represent any FOL formula.

\item \textbf{Generativity.} The coding scheme must define a joint
probability distribution over propositions---it must assign
$P(D \mid H)$ for any dataset $D$ and hypothesis $H$. Without
generativity, you cannot compute $L(D \mid H)$ and the MDL objective
is not defined. This rules out pure theorem provers, which give truth
values but not probabilities.

\item \textbf{Tractability.} The likelihoods $P(D \mid H)$ must be
computable in reasonable time. Exact inference in general Bayesian
networks is \#P-complete; a universal scheme needs a tractable
approximation.

\item \textbf{Structural transparency.} The coding scheme should make
the structure of a hypothesis explicit---which parts are the hard
logical constraints and which parts are the learned statistical
weights. This is what allows $L(H)$ to be decomposed into a structure
cost and a weight cost, connecting the coding scheme to the
Kolmogorov prior's notion of program length.

\end{enumerate}

No existing system satisfies all four. Pure FOL theorem provers
satisfy (1) and (4) but not (2) or (3). Markov Logic
Networks~\cite{coppola2024qbbn} satisfy (1) and (2) but not (3):
inference is \#P-complete and requires MCMC approximation. Neural
networks satisfy (2) and (3) but not (1) or (4): they approximate
logical structure without representing it explicitly, and their
hypothesis space is the space of weight configurations, not logical
propositions.

\subsection{The QBBN as the Best Known Candidate}

The Quantified Boolean Bayesian Network
(QBBN)~\cite{coppola2024qbbn} is the best known candidate for a
universal practical coding scheme. It satisfies all four properties.

\textbf{Universality (FOL completeness).} The QBBN paper proves that
any first-order logical formula can be expressed as a QBBN factor
graph. The QBBN is complete over FOL in the same sense that FOL is
complete over mathematics and science: any claim that can be stated
precisely can be stated as a QBBN.

\textbf{Generativity.} The QBBN defines a joint probability
distribution over boolean propositions via its factor graph. The AND
gates are deterministic (they implement exact boolean conjunction);
the OR gates are learned linear exponential models that assign
conditional probabilities. Together they define $P(D \mid H)$ for any
dataset $D$ of boolean observations.

\textbf{Tractability.} Inference in the QBBN uses belief propagation
(BP) on the AND/OR bipartite factor graph. Exact BP is polynomial on
trees; on loopy graphs it is an approximation, but the QBBN paper
reports empirical convergence in all tested cases. The bipartite
AND/OR structure---alternating conjunction and disjunction
layers---is the key architectural choice that makes BP tractable:
conjunction nodes are deterministic (no learned parameters, $O(2^n)$
but constant in practice for small conjuncts), and disjunction nodes
admit $O(n)$ updates under the Noisy-OR approximation.

\textbf{Structural transparency.} The AND/OR decomposition makes the
structure of a hypothesis explicit. AND gates represent the hard
logical constraints---the ``program'' part of the hypothesis. OR gates
represent the learned statistical weights---the ``parameter'' part.
The description length of the hypothesis is:
\[
  L(H) = L(\text{AND structure}) + L(\text{OR weights})
\]
The AND structure is the set of implication links---which premises
conjoin to produce which conclusions. Each link costs a fixed number
of bits to specify (its type and participating propositions). The OR
weights are the conditional probabilities learned for each disjunction
node; each weight costs $w$ bits to encode at $w$-bit precision. This
maps directly onto the Kolmogorov prior: AND structure corresponds to
program structure (fixed overhead per instruction), OR weights
correspond to program parameters (variable, proportional to
precision).

The \texttt{mdl-science} experiments in
Section~\ref{sec:experiments-logic} instantiate this coding scheme
directly. Each experiment uses QBBN theories---implication links and
conjunctive links---and computes MDL scores that reflect the true
$L(H) + L(D \mid H)$ under QBBN coding. The results (conjunctive
structure wins, theories are amortized, MDL selects correct
complexity) are consequences of the coding scheme's alignment with the
true generative structure.

\subsection{The Transformer as Independent Confirmation}

The coding scheme argument gains additional force from an independent
direction: the \texttt{shannon} paper~\cite{coppola2026} proves that
transformers implement exact belief propagation on QBBN-structured
factor graphs.

Specifically, the paper shows that:
\begin{itemize}
  \item Attention heads implement AND-type joint beliefs over
        query-key pairs---the AND gate computation of the QBBN.
  \item Feed-forward networks implement OR-type disjunctive
        updates---the OR gate computation of the QBBN.
  \item The residual stream is the shared message-passing bus---the
        variable nodes of the QBBN factor graph.
\end{itemize}

This means that a transformer trained by gradient descent on
next-token prediction is implicitly minimizing $L(D \mid H)$ where
$H$ is parameterized as a QBBN factor graph. The QBBN coding scheme
is not a design choice imposed on top of transformer training---it is
what transformer training discovers.

This has a direct implication for the Karpathy gap
(Section~\ref{sec:experiments-code}). The gap we demonstrated on toy
sort programs---full MDL achieves 65.2\% improvement while
$L(D \mid H)$-only achieves 0\%---is not a toy result. It applies
to every transformer trained without an explicit $L(H)$ penalty. A
language model trained to minimize validation loss (= $L(D \mid H)$
in QBBN coding) without any constraint on $L(H)$ will accumulate
complexity in the OR weights whenever doing so does not hurt
validation loss. The weight matrix grows without any signal that
something is wrong. Adding an explicit $L(H)$ penalty---compressed
size of the weight matrix, added to the training objective---would
complete the MDL objective in QBBN coding.

\subsection{The Open Problem}

The QBBN satisfies the four properties above for the domain of
logical/propositional hypotheses. It does not solve the general
coding scheme problem.

The open problem is cross-domain comparison. The QBBN gives a
principled MDL ranking of QBBN theories against each other. It does
not give a principled MDL ranking of a QBBN theory against, say, a
differential equation or a statistical regression model---because
there is no shared coding scheme that covers all three. Such a ranking
exists in principle: the Kolmogorov prior assigns description lengths
to all computable hypotheses, and the ranking is well-defined up to a
constant. But computing it requires $\K(H)$, which is uncomputable.

This is the computable residue of Hume's problem
(Section~\ref{sec:connections}). Hume showed that induction cannot be
deductively justified; we showed (in the Hume/Halting connection) that
this is equivalent to the uncomputability of $\K$. The coding scheme
problem is the same impossibility in engineering form: you cannot
build a universal computable coding scheme for all scientific
hypotheses for the same reason you cannot solve the halting problem.

What you can do---and what this paper does---is:
\begin{enumerate}
  \item Identify the best known candidate for a universal coding
        scheme within the logical/propositional domain (QBBN).
  \item Show that this coding scheme connects to the dominant
        practical architecture (transformers) via the
        \texttt{shannon} result.
  \item Demonstrate empirically that MDL in QBBN coding selects
        correct structure, penalizes spurious complexity, and
        amortizes over data as the theory predicts.
  \item Name the open problem honestly: a universal computable coding
        scheme covering all scientific domains does not yet exist.
\end{enumerate}

The Kolmogorov framework tells us such a scheme would exist if $\K$
were computable. It also tells us $\K$ is not computable. The QBBN is
the closest computable approximation currently known for the domain
where science and logic intersect---which, given FOL's expressive
completeness over mathematics and science, is most of the domain that
matters.